\chapter{Deep Q-Networks and Double DQN}

Chapter 7 put $Q$ into a parameterized approximator and showed the risk when function approximation, bootstrapping, and off-policy learning appear together---the \textbf{deadly triad}. Linear on-policy methods (SARSA FA) drop the off-policy piece and are stable under standard conditions.

In practice a richer approximator than linear FA is desirable---for example for images, raw sensors, or large state spaces. A neural $\hat{Q}(s,a;\theta)$ provides that, while control still keeps an off-policy Q-learning-style $\max$ backup. With bootstrapping, the full deadly triad is present again, so a neural off-policy semi-gradient update can diverge or train unstably, just as linear off-policy FA can. The stability result for on-policy linear TD and SARSA FA from Section 7.2 does not apply to this neural off-policy case.

The goal is still neural FA with an off-policy $\max$ backup; the question is how to train it without that instability. \textbf{DQN} (\citetxt{mnih2015}{Mnih et al., 2015}) answers with two changes to the basic neural Q-learning update: sample past transitions from an \textbf{experience replay} buffer, and bootstrap from a slowly updated \textbf{target network}. That is still ordinary DQN. A later refinement, \textbf{Double DQN} (\citetxt{vanhasselt2016}{van Hasselt et al., 2016}), does not replace those two ideas; it only changes how the $\max$ target is formed---online net chooses the greedy action, target net evaluates it---to reduce value overestimation.

Section 8.1 writes the neural update. Sections 8.2--8.3 cover replay and the target network; 8.4 the DQN loop and updates (8.3)--(8.6); 8.5 adds Double DQN.

\section{Neural action-values}

In linear FA, a fixed map produces a feature vector $\mathbf{x}(s,a)\in\mathbb{R}^d$ for each pair, and

\[
\hat{Q}(s,a;\mathbf{w})=\mathbf{w}^\top\mathbf{x}(s,a).
\]

Only $\mathbf{w}$ is learned. A neural network estimates the same $Q(s,a)$, but the features are no longer a fixed $\mathbf{x}$. They are computed inside the network from parameters $\theta$. Write the result as $\hat{Q}(s,a;\theta)$. When $\mathcal{A}$ is discrete and small, the network often maps the state $s$ to a vector of scores, one entry per action. Action choice is unchanged from Chapter 6: pick $A_t$ by $\varepsilon$-greedy or softmax on those scores.

The semi-gradient Q-learning update (7.3) lifts to

\begin{equation}
\theta
\leftarrow
\theta
+
\alpha
\bigl(
R_{t+1}
+
\gamma\max_{a'}\hat{Q}(S_{t+1},a';\theta)
-
\hat{Q}(S_t,A_t;\theta)
\bigr)
\nabla_\theta\hat{Q}(S_t,A_t;\theta).
\tag{8.1}
\end{equation}

The parameters $\theta$ play the same role as $\mathbf{w}$ in linear FA: they are the weights of $\hat{Q}$, updated by a semi-gradient step on the TD error. The difference is only the form of $\hat{Q}$---linear in $\mathbf{w}$, nonlinear in $\theta$---so the factor after the TD error is $\nabla_\theta\hat{Q}$ from backpropagation instead of $\mathbf{x}_t$.

\section{Experience replay}

Equation (8.1) updates $\theta$ from one transition $(S_t,A_t,R_{t+1},S_{t+1})$. Call that an \textbf{experience}. In the buffer we also store a done flag $d$: $d=1$ if $S_{t+1}$ is terminal (episode ended), and $d=0$ otherwise. So each stored experience is a tuple $(s,a,r,s',d)$.

In ordinary online TD-style learning, each new experience is used once in the update and then discarded. Successive experiences on a trajectory are strongly correlated, so successive gradient steps pull $\theta$ in similar directions. Rare experiences never return.

\textbf{Experience replay} does two things. First, \textbf{store} experiences in a buffer $\mathcal{D}$ instead of discarding them (a fixed-capacity queue of past $(s,a,r,s',d)$ tuples). Second, \textbf{replay} samples experiences from $\mathcal{D}$ after many steps (usually a random minibatch) and uses those samples in the learning update (8.1), rather than only the newest online transition.

Random replay mixes old and new experiences, so learning batches are less correlated than a raw online stream. The same experience can also be used in many updates. The backup is still $\max_{a'}$, so the method remains off-policy; replay only changes how data are stored and reused. Section 8.3 turns to a second issue: the bootstrap target in (8.1) still depends on the same $\theta$ being updated.

\section{Target network}

Replay (Section 8.2) changes \textbf{which} experiences enter the update: sample from $\mathcal{D}$ instead of using only the newest transition. It does not change \textbf{how} the one-step target is built. Inside (8.1) that target is still

\[
U_t
=
R_{t+1}
+
\gamma\max_{a'}\hat{Q}(S_{t+1},a';\theta)
\]

(with $U_t=R_{t+1}$ if $S_{t+1}$ is terminal). Both the prediction $\hat{Q}(S_t,A_t;\theta)$ and the bootstrap $\max_{a'}\hat{Q}(S_{t+1},a';\theta)$ depend on the same online weights $\theta$.

That shared dependence is unstable. Semi-gradient already treats $U_t$ as independent of $\theta$ when differentiating (Chapter 7), but after each update $\theta$ has changed, so the next target is recomputed from the new weights. With a nonlinear $\hat{Q}$, those successive targets can change abruptly: every step that improves the prediction also moves the bootstrap that defines the target. Replay does not fix this; it only supplies different experiences $(s,a,r,s',d)$ to the same update whose target still tracks $\theta$.

A \textbf{target network} removes the shared dependence. Keep a second parameter vector $\theta^-$ used only in the bootstrap $\max$. Update the online weights $\theta$ as usual; leave $\theta^-$ fixed for $C$ environment steps, then set $\theta^-\leftarrow\theta$ (or move $\theta^-$ slowly toward $\theta$). Between copies, replayed updates move $\hat{Q}(\cdot;\theta)$ toward a bootstrap that does not change with every gradient step. So $\theta^-$ is not a second learner: it is a periodically refreshed copy of $\theta$, introduced so the target stays fixed while $\theta$ is trained.

For a replayed experience $(s,a,r,s',d)$, replace $\theta$ by $\theta^-$ in the $\max$:

\begin{equation}
y
=
r
+
\gamma(1-d)\max_{a'}\hat{Q}(s',a';\theta^-).
\tag{8.2}
\end{equation}

When $d=1$ (terminal $s'$), the factor $(1-d)$ is $0$, so the $\max\hat{Q}$ term drops out and $y=r$: there is no bootstrap return from a terminal state. When $d=0$, (8.2) is $r+\gamma\max_{a'}\hat{Q}(s',a';\theta^-)$. So (8.2) is $U_t$ from (8.1), with $\theta^-$ in the bootstrap and $(1-d)$ writing the terminal case.

DQN samples a minibatch of $B$ experiences $\{(s_j,a_j,r_j,s'_j,d_j)\}_{j=1}^{B}$ from $\mathcal{D}$ and applies (8.2) to each:

\begin{equation}
y_j
=
r_j
+
\gamma(1-d_j)\max_{a'}\hat{Q}(s'_j,a';\theta^-),
\qquad
j=1,\ldots,B.
\tag{8.3}
\end{equation}

Neither side of that difference is the true $Q^*$: $y_j$ is a bootstrap TD target from (8.2), and $\hat{Q}(s_j,a_j;\theta)$ is the online prediction. The one-step error on sample $j$ is still written $y_j-\hat{Q}(s_j,a_j;\theta)$, treating $y_j$ as the regression target. Averaging the squared errors gives the minibatch objective

\begin{equation}
J(\theta)
=
\frac{1}{2B}
\sum_{j=1}^{B}
\bigl(y_j-\hat{Q}(s_j,a_j;\theta)\bigr)^2.
\tag{8.4}
\end{equation}

Each $y_j$ depends on $\theta^-$ (held fixed for the copy period), not on the online $\theta$ being differentiated. The \textbf{semi-gradient} of (8.4) therefore treats every $y_j$ as constant:

\begin{equation}
\nabla_\theta J(\theta)
=
-
\frac{1}{B}
\sum_{j=1}^{B}
\bigl(y_j-\hat{Q}(s_j,a_j;\theta)\bigr)
\nabla_\theta\hat{Q}(s_j,a_j;\theta).
\tag{8.5}
\end{equation}

Gradient descent on $J$ is then

\begin{equation}
\theta
\leftarrow
\theta
-
\alpha\,\nabla_\theta J(\theta)
=
\theta
+
\frac{\alpha}{B}
\sum_{j=1}^{B}
\bigl(y_j-\hat{Q}(s_j,a_j;\theta)\bigr)
\nabla_\theta\hat{Q}(s_j,a_j;\theta).
\tag{8.6}
\end{equation}

Compared with (8.1), the experiences come from replay and the bootstrap uses $\theta^-$ through the $y_j$. When $B=1$, (8.3)--(8.6) reduce to a single replayed transition.

\section{DQN}

\begin{implnote}[Algorithm (DQN)]
Input: network architecture for $\hat{Q}(s,a;\theta)$; replay capacity $N$; batch size $B$; discount $\gamma$; target-copy period $C$; step size $\alpha$ (or optimizer); exploration $\varepsilon$ (possibly scheduled); number of episodes $M$.

Output: online weights $\theta$; greedy policy $\pi(s)=\arg\max_a\hat{Q}(s,a;\theta)$.

\begin{enumerate}
\item Initialize online weights $\theta$ (e.g. randomly), set target weights to a copy $\theta^-\leftarrow\theta$, empty the replay buffer $\mathcal{D}$ (capacity $N$), and set the global step counter $k\leftarrow 0$.
\item Loop for episode $=1,2,\ldots,M$:
\begin{itemize}
\item Initialize environment and observe starting state $S$.
\item Loop until $S$ is terminal:
\begin{itemize}
\item Choose $A$ from $S$ using a policy based on $\hat{Q}(\cdot,\cdot;\theta)$ (e.g. $\varepsilon$-greedy).
\item Take $A$; observe $R$, $S'$, and done flag $d$ ($d=1$ if $S'$ terminal, else $0$).
\item Store experience $(S,A,R,S',d)$ in $\mathcal{D}$ (evict oldest if $|\mathcal{D}|=N$).
\item If $|\mathcal{D}|\geq B$:
\begin{itemize}
\item Sample a minibatch of $B$ experiences $\{(s_j,a_j,r_j,s'_j,d_j)\}_{j=1}^{B}$ uniformly from $\mathcal{D}$.
\item For each $j=1,\ldots,B$:
\begin{itemize}
\item $y_j\leftarrow r_j+\gamma(1-d_j)\max_{a'}\hat{Q}(s'_j,a';\theta^-)$.
\end{itemize}
\item $\theta\leftarrow\theta+\dfrac{\alpha}{B}\sum_{j=1}^{B}\bigl(y_j-\hat{Q}(s_j,a_j;\theta)\bigr)\nabla_\theta\hat{Q}(s_j,a_j;\theta)$.
\end{itemize}
\item $k\leftarrow k+1$.
\item If $k\bmod C=0$: set $\theta^-\leftarrow\theta$.
\item $S\leftarrow S'$.
\end{itemize}
\end{itemize}
\item Return $\theta$ and $\pi(s)=\arg\max_a\hat{Q}(s,a;\theta)$.
\end{enumerate}

\end{implnote}

\lstinputlisting{../code/ch08/dqn.py}


\section{Double DQN}

DQN already uses replay and a target network $\theta^-$. Its target is still a $\max$ over next-action values under $\theta^-$. That $\max$ both \textbf{selects} the greedy next action and \textbf{evaluates} it, using the same weights for both.

Because $\hat{Q}$ is noisy, $\arg\max$ returns the action with the highest current estimate---including when that estimate is too high by noise---and that value enters $y$. Learned values therefore tend to run \textbf{high} (\textbf{overestimation}). Replay and $\theta^-$ do not break this coupling.

That select-and-evaluate split is the Double Q-learning idea of \citetxt{vanhasselt2010}{van Hasselt, 2010}: one estimate picks the greedy action, another scores it, so a noisy $\max$ is not used for both jobs. \textbf{Double DQN} (\citetxt{vanhasselt2016}{van Hasselt et al., 2016}) keeps DQN's replay buffer $\mathcal{D}$ and target weights $\theta^-$ and applies the same split with the online net and the target net: $\theta$ selects the greedy action; $\theta^-$ evaluates it. A high score under $\theta$ need not be high under $\theta^-$, so it is not copied straight into $y$.

Ordinary DQN:

\[
y
=
r
+
\gamma(1-d)\max_{a'}\hat{Q}(s',a';\theta^-),
\]

i.e. select and evaluate with $\theta^-$:

\[
a^*
=
\arg\max_{a'}\hat{Q}(s',a';\theta^-),
\qquad
y
=
r
+
\gamma(1-d)\,\hat{Q}(s',a^*;\theta^-).
\]

Double DQN selects with $\theta$ instead:

\[
a^*
=
\arg\max_{a'}\hat{Q}(s',a';\theta),
\qquad
y
=
r
+
\gamma(1-d)\,\hat{Q}(s',a^*;\theta^-).
\]

For a minibatch:

\begin{equation}
y_j
=
r_j
+
\gamma(1-d_j)\,
\hat{Q}\Bigl(
s'_j,
\arg\max_{a'}\hat{Q}(s'_j,a';\theta);
\theta^-
\Bigr).
\tag{8.7}
\end{equation}

Replace (8.3) by (8.7); the loss (8.4)--(8.6), replay, and $\theta^-\leftarrow\theta$ stay unchanged.

In the program of Section 8.4, replace the \texttt{with torch.no\_grad():} block by

\lstinputlisting{../code/ch08/double_dqn_target.py}


DQN (and Double DQN) stay \textbf{value-based}. They learn $\hat{Q}$ and choose actions from it. Chapter 9 instead parameterizes $\pi_\theta(a\mid s)$ directly and optimizes expected return by policy gradients.
